% ============================================================
\chapter{Applications et \'ethique}
\label{ch:applications}
% ============================================================

\begin{intuition}
L'apprentissage profond n'est pas une fin en soi : c'est un outil au service de la science, de la m\'edecine, de l'industrie et de la soci\'et\'e. Ce chapitre explore les applications les plus marquantes et les questions \'ethiques fondamentales que tout praticien doit se poser.
\end{intuition}

% ------------------------------------------------------------
\section{Vision par ordinateur}
\label{sec:app_vision}
\index{vision par ordinateur}

\subsection{D\'etection d'objets}
\index{d\'etection d'objets}
\index{YOLO}

La d\'etection d'objets localise et classifie simultan\'ement plusieurs objets dans une image. L'approche YOLO (\emph{You Only Look Once}) divise l'image en une grille et pr\'edit directement les bo\^ites englobantes et les probabilit\'es de classe.

Pour chaque cellule $(i,j)$ de la grille, le mod\`ele pr\'edit :
\[
\hat{y}_{ij} = (p_{\mathrm{obj}}, b_x, b_y, b_w, b_h, c_1, \ldots, c_K)
\]
o\`u $p_{\mathrm{obj}}$ est la probabilit\'e de pr\'esence d'un objet, $(b_x, b_y, b_w, b_h)$ d\'efinissent la bo\^ite englobante, et $c_k$ sont les probabilit\'es de classe.

\subsection{Segmentation s\'emantique}
\index{segmentation}

La segmentation assigne une classe \`a chaque pixel. L'architecture U-Net\index{U-Net} (Ronneberger et al., 2015) utilise un encodeur-d\'ecodeur avec des connexions de saut :

\begin{center}
\begin{tikzpicture}[
    block/.style={rectangle, draw=NavyBlue, thick, minimum width=1.2cm, minimum height=0.6cm, fill=blue!10, font=\small},
    skip/.style={->, dashed, gray, thick},
    >=Stealth
]
\node[block] (e1) at (0,0) {Conv 64};
\node[block] (e2) at (1.5,-1) {Conv 128};
\node[block] (e3) at (3,-2) {Conv 256};
\node[block] (bn) at (4.5,-3) {Conv 512};
\node[block] (d3) at (6,-2) {Conv 256};
\node[block] (d2) at (7.5,-1) {Conv 128};
\node[block] (d1) at (9,0) {Conv 64};

\draw[->, thick] (e1) -- (e2);
\draw[->, thick] (e2) -- (e3);
\draw[->, thick] (e3) -- (bn);
\draw[->, thick] (bn) -- (d3);
\draw[->, thick] (d3) -- (d2);
\draw[->, thick] (d2) -- (d1);

\draw[skip] (e1) to[bend left=30] (d1);
\draw[skip] (e2) to[bend left=25] (d2);
\draw[skip] (e3) to[bend left=20] (d3);
\end{tikzpicture}
\end{center}

\subsection{G\'en\'eration d'images}
\index{g\'en\'eration d'images}

Les mod\`eles de diffusion (chapitre~\ref{ch:diffusion}) ont r\'evolutionn\'e la g\'en\'eration d'images. Stable Diffusion op\`ere dans l'espace latent d'un VAE (chapitre~\ref{ch:vae}), r\'eduisant consid\'erablement le co\^ut de calcul.

% ------------------------------------------------------------
\section{Traitement du langage naturel}
\label{sec:app_nlp}
\index{traitement du langage naturel}

\subsection{Mod\'elisation du langage}
\index{mod\`ele de langage}

Les grands mod\`eles de langage (LLM)\index{LLM} g\'en\`erent du texte de mani\`ere auto-r\'egressive :
\[
p(\bm{x}) = \prod_{t=1}^T p(x_t \mid x_1, \ldots, x_{t-1}; \theta)
\]

\begin{sota}
\begin{itemize}
\item \textbf{GPT-4} (OpenAI, 2023) : mod\`ele multimodal avec capacit\'es de raisonnement avanc\'ees.
\item \textbf{Claude} (Anthropic, 2024--2025) : accent sur la s\'ecurit\'e et l'alignement.
\item \textbf{LLaMA 3} (Meta, 2024) : mod\`ele ouvert performant.
\item \textbf{Gemini} (Google, 2024) : nativement multimodal.
\end{itemize}
\end{sota}

\subsection{Traduction automatique}

L'architecture Transformer\index{Transformer} (chapitre~\ref{ch:transformers}) a permis des avanc\'ees majeures en traduction, atteignant des performances quasi humaines sur de nombreuses paires de langues.

\subsection{R\'esum\'e automatique et r\'eponse aux questions}

Les mod\`eles pr\'e-entra\^in\'es comme BERT (bidirectionnel) et GPT (auto-r\'egressif) sont fine-tun\'es pour des t\^aches sp\'ecifiques de compr\'ehension.

% ------------------------------------------------------------
\section{Applications scientifiques}
\label{sec:app_science}
\index{applications scientifiques}

\subsection{Pr\'ediction de structures prot\'eiques}
\index{AlphaFold}

AlphaFold 2 (Jumper et al., 2021) a r\'esolu le probl\`eme du repliement des prot\'eines, vieux de 50 ans. L'architecture combine :
\begin{itemize}
\item Un Transformer avec attention \'evolutionnaire sur les alignements de s\'equences multiples (MSA).
\item Un module de structure qui pr\'edit les coordonn\'ees 3D.
\item Un raffinement it\'eratif des pr\'edictions.
\end{itemize}

\begin{remarque}
AlphaFold 3 (2024) \'etend la pr\'ediction aux complexes prot\'eine-ligand, prot\'eine-ADN et prot\'eine-ARN, unifiant la mod\'elisation des interactions mol\'eculaires.
\end{remarque}

\subsection{Pr\'evision m\'et\'eorologique}
\index{GraphCast}

GraphCast (Lam et al., 2023) utilise des r\'eseaux de neurones sur graphes pour la pr\'evision m\'et\'eorologique globale, surpassant les mod\`eles num\'eriques traditionnels tout en \'etant 1000$\times$ plus rapide.

\subsection{D\'ecouverte de m\'edicaments}

Les mod\`eles g\'en\'eratifs explorent l'espace chimique pour proposer de nouvelles mol\'ecules candidates :
\begin{itemize}
\item G\'en\'eration de mol\'ecules par VAE ou mod\`eles de diffusion.
\item Pr\'ediction d'affinit\'e mol\'eculaire par GNN (Graph Neural Networks).
\item Optimisation multi-objectifs (efficacit\'e, toxicit\'e, synth\'etisabilit\'e).
\end{itemize}

% ------------------------------------------------------------
\section{Chronologie des avanc\'ees majeures}

\begin{center}
\begin{tikzpicture}[
    event/.style={rectangle, draw=NavyBlue, thick, rounded corners, fill=blue!10, text width=3.5cm, font=\small, align=center},
    >=Stealth
]
\draw[thick, ->] (0,0) -- (14,0);

\foreach \x/\yr in {0.5/2012, 2.5/2014, 4.5/2015, 6.5/2017, 8.5/2020, 10.5/2022, 12.5/2024} {
    \draw (\x, -0.2) -- (\x, 0.2);
    \node[below, font=\small] at (\x, -0.3) {\yr};
}

\node[event, above] at (0.5, 0.5) {AlexNet\\ImageNet};
\node[event, below] at (2.5, -0.8) {GAN\\VAE};
\node[event, above] at (4.5, 0.5) {ResNet\\BatchNorm};
\node[event, below] at (6.5, -0.8) {Transformer\\Attention};
\node[event, above] at (8.5, 0.5) {GPT-3\\Diffusion};
\node[event, below] at (10.5, -0.8) {ChatGPT\\Stable Diff.};
\node[event, above] at (12.5, 0.5) {GPT-4\\Claude\\Sora};
\end{tikzpicture}
\end{center}

% ------------------------------------------------------------
\section{Interpr\'etabilit\'e}
\label{sec:interpretability}
\index{interpr\'etabilit\'e}

\subsection{Cartes de saillance}
\index{carte de saillance}

Les cartes de saillance identifient les r\'egions de l'entr\'ee les plus influentes pour la pr\'ediction :
\[
S_i = \left|\frac{\partial f_\theta(\bm{x})}{\partial x_i}\right|
\]

\subsection{Grad-CAM}
\index{Grad-CAM}

Grad-CAM (Selvaraju et al., 2017) utilise les gradients de la couche convolutive cible pour produire une carte de chaleur :
\[
\alpha_k^c = \frac{1}{Z}\sum_i \sum_j \frac{\partial y^c}{\partial A^k_{ij}}, \qquad L_{\mathrm{Grad\text{-}CAM}}^c = \relu\!\left(\sum_k \alpha_k^c A^k\right)
\]

\begin{codeblock}[title=\textbf{Grad-CAM en PyTorch}]
\begin{minted}{python}
import torch
import torch.nn.functional as F

class GradCAM:
    def __init__(self, model, target_layer):
        self.model = model
        self.gradients = None
        self.activations = None
        target_layer.register_forward_hook(
            lambda m, i, o: setattr(self, 'activations', o.detach()))
        target_layer.register_full_backward_hook(
            lambda m, gi, go: setattr(self, 'gradients', go[0].detach()))

    def __call__(self, x, class_idx=None):
        output = self.model(x)
        if class_idx is None:
            class_idx = output.argmax(1)
        self.model.zero_grad()
        one_hot = torch.zeros_like(output)
        one_hot[0, class_idx] = 1
        output.backward(gradient=one_hot)

        weights = self.gradients.mean(dim=[2, 3], keepdim=True)
        cam = F.relu((weights * self.activations).sum(1, keepdim=True))
        cam = F.interpolate(cam, size=x.shape[2:],
                           mode='bilinear', align_corners=False)
        cam = cam / (cam.max() + 1e-8)
        return cam.squeeze().detach().numpy()
\end{minted}
\end{codeblock}

\subsection{SHAP et valeurs de Shapley}
\index{SHAP}

Les valeurs de Shapley, issues de la th\'eorie des jeux coop\'eratifs, attribuent \`a chaque variable une contribution \'equitable \`a la pr\'ediction :
\[
\phi_i = \sum_{S \subseteq \{1,\ldots,d\} \setminus \{i\}} \frac{|S|!(d-|S|-1)!}{d!}\left[f(S \cup \{i\}) - f(S)\right]
\]

% ------------------------------------------------------------
\section{Biais et \'equit\'e}
\label{sec:fairness}
\index{biais}
\index{\'equit\'e}

\subsection{Sources de biais}

\begin{attention}
Les mod\`eles d'apprentissage profond peuvent amplifier les biais pr\'esents dans les donn\'ees d'entra\^inement :
\begin{itemize}
\item \textbf{Biais de repr\'esentation} : certains groupes sous-repr\'esent\'es dans les donn\'ees.
\item \textbf{Biais d'\'etiquetage} : annotations refl\'etant des pr\'ejug\'es humains.
\item \textbf{Biais d'agr\'egation} : un mod\`ele unique pour des populations h\'et\'erog\`enes.
\end{itemize}
\end{attention}

\subsection{M\'etriques d'\'equit\'e}

\begin{definition}[Parit\'e d\'emographique]
Un classificateur satisfait la parit\'e d\'emographique si :
\[
P(\hat{Y} = 1 \mid A = a) = P(\hat{Y} = 1 \mid A = b) \quad \forall a, b
\]
o\`u $A$ est l'attribut sensible.
\end{definition}

\begin{definition}[\'Egalit\'e des chances]
Un classificateur satisfait l'\'egalit\'e des chances si :
\[
P(\hat{Y} = 1 \mid Y = 1, A = a) = P(\hat{Y} = 1 \mid Y = 1, A = b) \quad \forall a, b
\]
\end{definition}

% ------------------------------------------------------------
\section{Confidentialit\'e}
\label{sec:privacy}
\index{confidentialit\'e}
\index{confidentialit\'e diff\'erentielle}

\subsection{Confidentialit\'e diff\'erentielle}

\begin{definition}[$(\varepsilon,\delta)$-confidentialit\'e diff\'erentielle]
Un m\'ecanisme al\'eatoire $\mathcal{M}$ satisfait la $(\varepsilon,\delta)$-confidentialit\'e diff\'erentielle si, pour tous ensembles de donn\'ees adjacents $D, D'$ (diff\'erant d'un seul exemple) et tout ensemble mesurable $S$ :
\[
P(\mathcal{M}(D) \in S) \leq e^\varepsilon \cdot P(\mathcal{M}(D') \in S) + \delta
\]
\end{definition}

DP-SGD (Abadi et al., 2016) applique la confidentialit\'e diff\'erentielle \`a la SGD en \'ecr\^etant les gradients individuels et en ajoutant du bruit gaussien :
\[
\tilde{\bm{g}}_t = \frac{1}{B}\left(\sum_{i \in \mathcal{B}} \mathrm{clip}(\bm{g}_i, C) + \mathcal{N}(0, \sigma^2 C^2 I)\right)
\]

\subsection{Apprentissage f\'ed\'er\'e}
\index{apprentissage f\'ed\'er\'e}

L'apprentissage f\'ed\'er\'e entra\^ine un mod\`ele global sans centraliser les donn\'ees. Chaque client $k$ calcule une mise \`a jour locale, et le serveur agr\`ege :
\[
\theta_{t+1} = \sum_{k=1}^K \frac{n_k}{n} \theta_{t+1}^{(k)}
\]

% ------------------------------------------------------------
\section{Impact environnemental}
\label{sec:environment}
\index{impact environnemental}

L'entra\^inement de grands mod\`eles consomme des ressources consid\'erables :

\begin{center}
\begin{tabular}{lrr}
\toprule
\textbf{Mod\`ele} & \textbf{Param\`etres} & \textbf{CO$_2$ estim\'e (tonnes)} \\
\midrule
BERT-Large & 340M & $\sim$0.6 \\
GPT-3 & 175B & $\sim$500 \\
LLaMA 2 70B & 70B & $\sim$290 \\
GPT-4 (estim\'e) & $>$1T & $\sim$5\,000+ \\
\bottomrule
\end{tabular}
\end{center}

\begin{bonnepratique}
\begin{itemize}
\item Privil\'egier le \textbf{fine-tuning} de mod\`eles pr\'e-entra\^in\'es plut\^ot que l'entra\^inement de z\'ero.
\item Utiliser des techniques d'efficacit\'e : quantification, distillation, LoRA.
\item Rapporter l'empreinte carbone dans les publications.
\item Choisir des centres de calcul aliment\'es en \'energies renouvelables.
\end{itemize}
\end{bonnepratique}

% ------------------------------------------------------------
\section{Cadre r\'eglementaire}
\label{sec:regulation}
\index{r\'eglementation}
\index{EU AI Act}

Le \textbf{EU AI Act} (2024) classe les syst\`emes d'IA par niveau de risque :
\begin{enumerate}
\item \textbf{Risque inacceptable} : interdit (notation sociale, manipulation subliminale).
\item \textbf{Haut risque} : r\'eglement\'e (sant\'e, justice, recrutement).
\item \textbf{Risque limit\'e} : obligations de transparence.
\item \textbf{Risque minimal} : libre utilisation.
\end{enumerate}

% ------------------------------------------------------------
\section{Mod\`eles multimodaux et perspectives}
\label{sec:multimodal}
\index{multimodal}

\begin{sota}
\begin{itemize}
\item \textbf{Mod\`eles multimodaux} : GPT-4V, Claude 3.5 Vision, Gemini combinent texte, image, audio et vid\'eo dans un seul mod\`ele.
\item \textbf{Agents IA} : syst\`emes autonomes utilisant des outils (navigation web, ex\'ecution de code, appels API).
\item \textbf{Mod\`eles du monde} : apprentissage de repr\'esentations physiques pour la robotique et la simulation (Sora, World Labs).
\item \textbf{IA scientifique} : acc\'el\'eration de la d\'ecouverte dans tous les domaines (mat\'eriaux, climat, biologie).
\end{itemize}
\end{sota}

% ------------------------------------------------------------
\section{R\'esum\'e du chapitre}

\begin{formules}
\begin{itemize}
\item \textbf{D\'etection} : YOLO pr\'edit bo\^ites englobantes et classes par cellule de grille
\item \textbf{Segmentation} : U-Net avec encodeur-d\'ecodeur et connexions de saut
\item \textbf{Grad-CAM} : $L^c = \relu(\sum_k \alpha_k^c A^k)$ avec $\alpha_k^c = \frac{1}{Z}\sum_{ij}\frac{\partial y^c}{\partial A^k_{ij}}$
\item \textbf{SHAP} : valeurs de Shapley pour l'attribution
\item \textbf{DP-SGD} : SGD avec \'ecr\^etage des gradients et bruit gaussien
\item \textbf{Apprentissage f\'ed\'er\'e} : agr\'egation pond\'er\'ee des mises \`a jour locales
\end{itemize}
\end{formules}

% ------------------------------------------------------------
\section{Exercices}

\begin{exercice}[$\star$ --- Analyse]
Montrer que la parit\'e d\'emographique et l'\'egalit\'e des chances ne peuvent pas \^etre satisfaites simultan\'ement sauf dans des cas d\'eg\'en\'er\'es. Quelles sont les implications pratiques de ce r\'esultat ?
\end{exercice}

\begin{exercice}[$\star$ --- Calcul]
Calculer le nombre de FLOPS n\'ecessaires pour un forward pass d'un Transformer avec $L=32$ couches, $d=4096$, $h=32$ t\^etes et une s\'equence de longueur $n=2048$. Estimer le co\^ut \'energ\'etique sur un GPU A100.
\end{exercice}

\begin{exercice}[$\star\star$ --- Impl\'ementation]
Impl\'ementer Grad-CAM pour un ResNet-18 pr\'e-entra\^in\'e sur ImageNet. Visualiser les cartes d'attention pour diff\'erentes classes et analyser les r\'esultats.
\end{exercice}

\begin{exercice}[$\star\star$ --- Impl\'ementation]
Impl\'ementer un pipeline d'apprentissage f\'ed\'er\'e simple avec 5 clients sur MNIST. Comparer la performance avec un entra\^inement centralis\'e et analyser l'impact du nombre de rounds de communication.
\end{exercice}

\begin{exercice}[$\star\star\star$ --- Projet de recherche]
Conduire un audit d'\'equit\'e complet sur un mod\`ele de classification (par exemple, pr\'ediction de revenu sur le dataset Adult). Calculer la parit\'e d\'emographique, l'\'egalit\'e des chances et l'\'egalit\'e pr\'edictive. Proposer et impl\'ementer une m\'ethode de d\'ebiaisage (pr\'e-traitement, en entra\^inement, ou post-traitement) et \'evaluer son impact sur la pr\'ecision et l'\'equit\'e.
\end{exercice}
